\documentclass[11pt,a4paper]{article}
\usepackage[utf8]{inputenc}
\usepackage[vietnamese]{babel}
\usepackage{amsmath,amssymb}
\usepackage{geometry}
\geometry{a4paper, total={180mm,260mm}, left=15mm, right=15mm, top=18mm, bottom=18mm}

\usepackage{tikz}
\usetikzlibrary{shapes.geometric, arrows.meta, positioning, calc}
\tikzset{
  arr/.style={-{Stealth[scale=0.8]}},
  every node/.append style={align=center}
}
\usepackage{tcolorbox}
\tcbuselibrary{skins,breakable}
\usepackage{fontawesome}
\usepackage{fancyhdr}
\usepackage{enumitem}

\pagestyle{fancy}
\fancyhf{}
\renewcommand{\headrulewidth}{0.5pt}
\renewcommand{\footrulewidth}{0.5pt}
\fancyhead[L]{\small\bfseries\color{violet!85!black}\faBook\ SỔ TAY KIẾN THỨC TOÁN 2}
\fancyhead[R]{\small\bfseries\color{violet!85!black}Lớp Toán Thầy Lĩnh -- loptoanthaylinh.vn}
\fancyfoot[C]{\small\color{gray}-- Sổ tay kiến thức: Bài 37. Phép nhân --}

\setlength{\parindent}{0pt}
\setlength{\parskip}{0pt}

\begin{document}

% --- BANNER BÀI HỌC ---
\begin{tcolorbox}[
  enhanced,
  colback=violet!8!white,
  colframe=violet!70!black,
  arc=3mm, boxrule=1.2pt,
  left=3mm, right=3mm, top=2mm, bottom=2mm,
  halign=center
]
  {\LARGE\bfseries\color{violet!90!black}\faBookmark\ BÀI 37: PHÉP NHÂN}\\[1mm]
  {\itshape\footnotesize\color{black!75} Hình thành khái niệm phép nhân từ tổng các số hạng bằng nhau trên 1 trang A4}
\end{tcolorbox}

\vspace{-2mm}

\noindent
\begin{minipage}[t]{0.49\textwidth}\vspace{0pt}
\begin{tcolorbox}[
  enhanced,
  colback=violet!4!white,
  colframe=violet!75!black,
  arc=2mm, boxrule=0.8pt,
  left=2mm, right=2mm, top=1.8mm, bottom=1.8mm,
  title={\footnotesize\bfseries\color{white}\faKey\ 1. Khái niệm phép nhân}
]
{\footnotesize
\begin{itemize}\setlength{\itemsep}{1pt}\setlength{\leftskip}{-2.5mm}
  \item Phép nhân là cách viết gọn của \textbf{tổng các số hạng bằng nhau}.
  \item \textbf{Ví dụ:} Mỗi đĩa có $2$ quả cam, có $3$ đĩa như thế.
  $$\text{Tổng: } 2 + 2 + 2 = 6$$
  Số $2$ được lấy $3$ lần, ta viết thành phép nhân:
  $$\mathbf{2 \times 3 = 6}$$
  \item Dấu $\times$ đọc là \textbf{"nhân"}. Phép tính đọc là: \textit{"Hai nhân ba bằng sáu"}.
\end{itemize}
}
\end{tcolorbox}
\end{minipage}
\hfill
\begin{minipage}[t]{0.49\textwidth}\vspace{0pt}
\begin{tcolorbox}[
  enhanced,
  colback=teal!4!white,
  colframe=teal!75!black,
  arc=2mm, boxrule=0.8pt,
  left=2mm, right=2mm, top=1.8mm, bottom=1.8mm,
  title={\footnotesize\bfseries\color{white}\faSliders\ 2. Kĩ năng chuyển đổi tổng $\leftrightarrow$ tích}
]
{\footnotesize
\begin{itemize}\setlength{\itemsep}{1pt}\setlength{\leftskip}{-2.5mm}
  \item \textbf{Chuyển từ tổng sang nhân:}
  $$a + a + \dots + a \text{ ($n$ số hạng $a$)} = a \times n$$
  \textit{Ví dụ:} $5 + 5 + 5 + 5 = 5 \times 4 = 20$.
  \item \textbf{Chuyển từ nhân sang tổng:}
  $$a \times n = \underbrace{a + a + \dots + a}_{n\text{ số hạng } a}$$
  \textit{Ví dụ:} $3 \times 4 = 3 + 3 + 3 + 3 = 12$.
  \item Lưu ý: Số đứng trước là \textbf{giá trị mỗi phần}, số đứng sau là \textbf{số lần lấy}.
\end{itemize}
}
\end{tcolorbox}
\end{minipage}

\vspace{-1.5mm}
\begin{tcolorbox}[
  enhanced,
  colback=orange!4!white,
  colframe=orange!80!black,
  arc=2mm, boxrule=0.8pt,
  left=2mm, right=2mm, top=1.5mm, bottom=1.5mm,
  title={\footnotesize\bfseries\color{white}\faSitemap\ 3. Sơ đồ mô hình trực quan: Từ phép cộng sang phép nhân}
]
\centering
\begin{tikzpicture}[x=1.5cm, y=0.8cm]
  \node[rectangle, rounded corners=2pt, draw=violet!80!black, fill=violet!15, font=\scriptsize\bfseries] (c1) at (0,0) {ĐĨA 1\\$2$ quả};
  \node[font=\normalsize\bfseries] at (1,0) {$+$};
  \node[rectangle, rounded corners=2pt, draw=violet!80!black, fill=violet!15, font=\scriptsize\bfseries] (c2) at (2,0) {ĐĨA 2\\$2$ quả};
  \node[font=\normalsize\bfseries] at (3,0) {$+$};
  \node[rectangle, rounded corners=2pt, draw=violet!80!black, fill=violet!15, font=\scriptsize\bfseries] (c3) at (4,0) {ĐĨA 3\\$2$ quả};

  \draw[->, very thick, red!85!black] (2,-0.6) -- node[right, font=\tiny\bfseries] {Viết gọn thành} (2,-1.3);

  \node[rectangle, rounded corners=2pt, draw=teal!80!black, fill=teal!15, font=\footnotesize\bfseries] at (2,-1.7) {
    $2$ được lấy $3$ lần $\implies \mathbf{2 \times 3 = 6}$
  };
\end{tikzpicture}
\end{tcolorbox}

\vspace{-1.5mm}
\noindent
\begin{minipage}[t]{0.49\textwidth}\vspace{0pt}
\begin{tcolorbox}[
  enhanced,
  colback=green!4!white,
  colframe=green!60!black,
  arc=2mm, boxrule=0.8pt,
  left=2mm, right=2mm, top=1.8mm, bottom=1.8mm,
  title={\footnotesize\bfseries\color{white}\faCheckCircle\ 4. Ví dụ thực hành chuẩn mực}
]
{\footnotesize
\textbf{1. Chuyển tổng thành phép nhân:}\\
- $4 + 4 + 4 + 4 + 4 = 4 \times 5 = \mathbf{20}$.\\
- $7 + 7 + 7 = 7 \times 3 = \mathbf{21}$.
\textbf{2. Tính kết quả phép nhân dựa vào tổng:}\\
Tính $6 \times 3$?\\
\textit{Giải:} $6 \times 3 = 6 + 6 + 6 = \mathbf{18}$.
}
\end{tcolorbox}
\end{minipage}
\hfill
\begin{minipage}[t]{0.49\textwidth}\vspace{0pt}
\begin{tcolorbox}[
  enhanced,
  colback=red!4!white,
  colframe=red!75!black,
  arc=2mm, boxrule=0.8pt,
  left=2mm, right=2mm, top=1.8mm, bottom=1.8mm,
  title={\footnotesize\bfseries\color{white}\faExclamationTriangle\ 5. Mẹo tránh bẫy trắc nghiệm}
]
{\footnotesize
\begin{itemize}\setlength{\itemsep}{1pt}\setlength{\leftskip}{-2.5mm}
  \item \textbf{Bẫy tổng các số không bằng nhau:}\\
  $2 + 3 + 4$ \textbf{không thể} viết thành phép nhân vì các số hạng không bằng nhau!
  \item \textbf{Bẫy viết ngược thừa số:}\\
  $2 + 2 + 2$ là $2$ lấy $3$ lần $\implies 2 \times 3$ (chuẩn bản chất). Tránh viết nhầm thành $3 \times 2$.
\end{itemize}
}
\end{tcolorbox}
\end{minipage}


\end{document}
